\newif\ifacmart
\IfFileExists{acmart.cls}{
  \documentclass[sigconf,nonacm]{acmart}
  \acmarttrue
}{
  \documentclass[10pt,twocolumn]{article}
  \acmartfalse
}

\usepackage{booktabs}
\usepackage{hyperref}
\ifacmart\else
  \usepackage[margin=0.75in]{geometry}
\fi

\title{Math Reasoning with Adaptive MCTS and Process-Level Scoring}

\ifacmart
  \author{Ruochen Wu}
  \affiliation{\institution{Columbia University}}
  \email{rw3152@columbia.edu}
  \author{Hammer Niu}
  \affiliation{\institution{Columbia University}}
  \email{hn2477@columbia.edu}
  \author{Shengbo Yi}
  \affiliation{\institution{Columbia University}}
  \email{sy3302@columbia.edu}
\else
  \author{Ruochen Wu (rw3152) \and Hammer Niu (hn2477) \and Shengbo Yi (sy3302)}
\fi

\begin{document}
\maketitle

\begin{abstract}
This project investigates whether search and process-level scoring can make
LLM-based math reasoning more reliable than a single direct reasoning path. We
implement a baseline Monte Carlo Tree Search (MCTS) solver, an adaptive search
variant, and a PPM-compatible process scorer that combines a learned Process
Preference Model with a deterministic verifier. The system is exposed through
command-line scripts, a Streamlit comparison interface, and a FastAPI backend.
Deterministic ablations show that adaptive search improves controlled answer
selection from 0/3 to 3/3 while reducing model calls from 6.0 to 2.0 per
problem, and that the hybrid scorer corrects 3/3 deliberately constructed
process-score mis-rankings. The repository also includes scripts for MATH Level
5 and OlympiadBench evaluation when API keys are available.
\end{abstract}

\section{Introduction and Problem Formulation}

Large language models can solve many short math problems, but direct generation
often commits to a single reasoning path. For multi-step problems, a small
intermediate error can propagate into a fluent but incorrect final answer. This
project studies a more deliberate alternative: search over multiple candidate
reasoning steps, score the quality of intermediate process steps, and return the
best trajectory.

The project goal is to reproduce a working math-reasoning baseline and improve
it with two concrete innovations: an adaptive MCTS search policy and a
process-level scorer that combines learned PPM scoring with a lightweight math
verifier. The central research question is: given the same math problem and the
same generator model, can structured search and process scoring improve which
reasoning path is selected?

The system is packaged as a runnable command-line, Streamlit, and FastAPI demo.
The final report, repository, and demo artifacts are intentionally aligned with
the course rubric so that the code, presentation, and evaluation describe the
same end-to-end system.

\noindent Repository:
\url{https://github.com/HammerNiu/Math-Reasoning}

\noindent Demo artifact:
\url{https://github.com/HammerNiu/Math-Reasoning/blob/main/docs/demo_video.gif}

\section{Related Work and Background}

The project is inspired by search-based reasoning systems such as rStar-Math,
where candidate reasoning trajectories are generated, scored, and improved over
multiple rounds. MCTS provides a natural framework for balancing exploration and
exploitation: selection chooses promising nodes, expansion adds candidate
reasoning steps, simulation estimates path quality, and backpropagation updates
node values. Process preference models complement this by scoring the quality of
intermediate steps rather than judging only the final answer.

This project also relates to verifier-guided generation. Final-answer
verification is often too late: a path can be fluent and internally
inconsistent before the answer line appears. Process supervision instead asks
whether each intermediate step is useful, mathematically grounded, and likely to
lead to a correct final state. Our implementation uses this idea in two forms:
a neural PPM trained on preference pairs and a deterministic verifier that
provides a fallback score when no trained checkpoint is available.

\section{System Design and Methodology}

The system has four main components. First, a model interface supports OpenAI,
DeepSeek, Anthropic, and local Ollama models. Second, the MCTS controller
generates and searches over candidate steps. Third, the process scorer ranks
steps before expansion. Fourth, the demo layer presents baseline and improved
systems side by side.

\subsection{End-to-End Workflow}

At inference time, the system receives a problem string and initializes it as
the root state of a search tree. The generator model proposes candidate next
steps using a prompt that asks for explicit \texttt{STEP:} lines and a
\texttt{FINAL ANSWER:} marker when the problem is solved. MCTS then selects and
expands nodes. For the improved system, candidate steps are de-duplicated,
scored, and pruned before expansion. After simulation, rewards are
backpropagated through the tree and the best child is returned as the selected
reasoning action.

The design keeps policy and reward responsibilities separate. The policy model
generates candidate steps. The reward side can use a trained PPM, a separate
critic model, the deterministic verifier, or a hybrid PPM-verifier scorer. This
separation makes the system easier to inspect and cheaper to adapt: improving
the scorer does not require changing the generator.

\subsection{Search Innovation}

The baseline MCTS preserves a simple search strategy. The adaptive strategy adds
larger branching near the root, duplicate removal, diversity-aware scoring,
weak-branch pruning, early stopping, evaluation caching, and context/token
controls. These changes are implemented in \path{src/core/mcts.py}.

The adaptive strategy changes three decisions. First, it asks for more
candidate steps near the root, where early commitment is most costly. Second, it
uses a simple value-diversity score to avoid spending all simulations on
near-duplicate branches. Third, it stops early when a high-confidence terminal
answer is found. These changes directly target the baseline failure mode:
following the first plausible but weak branch.

\subsection{Scoring Innovation}

The scoring-side innovation is a PPM-compatible verifier and hybrid scorer in
\path{src/core/scoring.py}. The verifier rewards mathematical progress,
explicit operations, verification language, and final-answer markers. It
penalizes vague guesses and unrelated steps. When a trained PPM checkpoint is
available, the hybrid scorer combines the learned PPM score with the verifier;
when no checkpoint is available, the verifier still provides a deterministic
process score.

The verifier is not intended to replace proof checking. Instead, it is a
low-cost process prior. It increases scores for steps that contain mathematical
operations, explicit solving verbs, or verification language, and decreases
scores for guesses, vague language, or unrelated actions. The hybrid scorer
normalizes raw PPM logits into a 0--1 range and interpolates them with the
verifier score. This gives the project a scoring-side innovation that works both
before and after PPM training.

\subsection{Demo and API Design}

The Streamlit interface is designed for the presentation rubric: it runs the
same input through baseline MCTS and improved MCTS, then displays selected
steps, reasoning trajectories, model-call counts, pruned branches, and latency.
The FastAPI backend exposes the same behavior through \texttt{/solve},
\texttt{/compare-models}, and \texttt{/health}. This makes the system usable as
both a live demo and a reproducible service endpoint.

\section{Implementation and Data}

The main implementation files are:

\begin{itemize}
    \item \path{src/core/mcts.py}: MCTS, adaptive search, top-k process pruning.
    \item \path{src/core/ppm.py}: neural process preference model and trainer.
    \item \path{src/core/scoring.py}: verifier and hybrid process scorer.
    \item \path{app.py}: side-by-side comparison demo.
    \item \path{backend/main.py}: FastAPI endpoints for solving and comparison.
    \item \path{tools/}: trajectory collection, PPM training, deterministic ablations.
    \item \path{experiments/run_experiment.py} and \path{eval.py}: MATH and OlympiadBench evaluation.
\end{itemize}

Trajectory data can be generated by running
\path{tools/collect_trajectories.py}. Preference pairs are then used to train
the PPM with \path{tools/train_ppm.py}. The full benchmark path uses MATH Level
5 and OlympiadBench through HuggingFace \texttt{datasets}.

\subsection{Preference Pair Construction}

The PPM training pipeline converts trajectories into pairwise preferences.
Steps from high-reward trajectories become preferred examples, while steps from
low-reward trajectories become non-preferred examples. The trainer embeds each
step using a local sentence-transformer model and optimizes a hinge-style
ranking loss so preferred steps receive higher values than non-preferred steps.
This design avoids fine-tuning the generator model and keeps training local.

\subsection{Configuration and Reproducibility}

Configuration lives in \path{config/default.json}. The repository includes both
cloud and local model paths: OpenAI, DeepSeek, Anthropic, and Ollama.
Deterministic ablation scripts require no API keys, while full benchmark runs
require model credentials and dataset access. The README lists exact commands
for compiling, running ablations, launching the demo, training PPM checkpoints,
and running benchmark comparisons.

\section{Evaluation and Results}

The committed evaluation contains deterministic no-key ablations plus scripts
for API-backed benchmark runs. The deterministic experiments are intentionally
small but isolate the two required algorithmic changes without depending on
cloud API randomness.

\subsection{Search Ablation}

The search ablation compares baseline MCTS against adaptive MCTS on three
scripted math problems. The baseline is intentionally exposed to a distracting
first candidate. Adaptive MCTS should identify the useful final step sooner.

\begin{table}[h]
\centering
\caption{Member 1 deterministic search ablation.}
\begin{tabular}{lrrrr}
\toprule
Strategy & Accuracy & Nodes & Calls & Pruned \\
\midrule
Baseline MCTS & 0/3 & 3.0 & 6.0 & 0.0 \\
Adaptive MCTS & 3/3 & 1.0 & 2.0 & 2.0 \\
\bottomrule
\end{tabular}
\end{table}

Qualitatively, the baseline chooses the first distracting candidate in each
controlled problem: an unchecked arithmetic or algebraic guess. Adaptive MCTS
promotes the terminal or mathematically grounded candidate because its
heuristics reward explicit solution progress and penalize vague guessing. This
supports the Member 1 claim that search-side pruning and early stopping improve
both quality and efficiency in controlled cases.

\subsection{Scoring Ablation}

The scoring ablation uses a deliberately weak scorer that overvalues guesses.
The hybrid verifier scorer corrects all three mis-rankings by preferring
mathematically useful steps.

\begin{table}[h]
\centering
\caption{Member 2 deterministic scoring ablation.}
\begin{tabular}{lrr}
\toprule
Task & Old scorer fixed? & Hybrid fixed? \\
\midrule
Linear equation & No & Yes \\
Addition & No & Yes \\
Derivative & No & Yes \\
\bottomrule
\end{tabular}
\end{table}

The scoring ablation is intentionally adversarial: the old scorer is biased
toward guesses. The hybrid scorer fixes this by applying the verifier prior.
For example, in the linear-equation case, the weak scorer prefers ``Guess
\(x=4\) without checking the equation,'' while the hybrid scorer prefers
``Subtract 3 from both sides to get \(x=2\).'' This demonstrates the task-split
requirement for error-analysis examples showing corrected mis-ranking cases.

\subsection{End-to-End Evaluation}

The Streamlit demo runs the same input problem through baseline MCTS and
improved MCTS with process scoring. It reports selected step, trajectory, model
calls, pruned branches, and latency. The full MATH/OlympiadBench protocol is
implemented in \path{experiments/run_experiment.py} and \path{eval.py}.

The full benchmark protocol compares three configurations on held-out MATH
Level 5 problems: baseline MCTS, adaptive MCTS, and adaptive MCTS with PPM
top-k pruning. Metrics include accuracy, average model calls, latency, and PPM
pruned branches. A separate evaluation script supports direct prompting,
MCTS-only, and MCTS+PPM on MATH or OlympiadBench. These scripts are included so
the final API-backed results can be reproduced exactly from the repository.

\subsection{Failure Analysis}

The main remaining risks are scorer noise and benchmark variance. A weak PPM
can prune useful branches if used without enough training data; therefore the
system supports verifier-only scoring and documents that top-k PPM pruning
should be enabled only after trajectory collection and training. Answer checking
for competition math is also imperfect because equivalent expressions can appear
in different formats. The evaluation scripts use boxed-answer extraction, LaTeX
normalization, numeric comparison, and token matching, but manual review remains
useful for borderline cases.

\section{Discussion}

The controlled ablations show that the two algorithmic modifications work as
intended: adaptive search reduces wasted exploration and the hybrid scorer
corrects process-level mis-rankings. The main limitation is that deterministic
ablations are smaller than full benchmark evaluation. The repository therefore
includes the complete benchmark scripts and instructions so the final API-backed
run can be reproduced before submission.

\subsection{Division of Work}

The implementation matches the three-person task split. Member 1 owns the MCTS
search modification and ablation. Member 2 owns the PPM/verifier scoring
modification and mis-ranking analysis. Member 3 owns the comparison demo,
repository organization, and evaluation harness. All three components are
represented in code, documentation, and presentation materials rather than being
separated into only coding or writing tasks.

\subsection{Submission Readiness}

The repository now contains the report source, rubric checklist, demo guide,
experiment analysis, deterministic result tables, and a demo GIF artifact. For
CourseWorks, the final step is to export this report as
\texttt{rw3152\_hn2477\_sy3302.pdf}, confirm the GitHub and demo links are
accessible, and optionally replace the GIF with a narrated hosted video.

Future work includes training a larger PPM on more trajectories, adding richer
mathematical verification, and evaluating on broader reasoning tasks beyond
math.

\section{Conclusion}

This project delivers a coherent math-reasoning system with a baseline, two
implemented innovations, runnable demos, deterministic evidence, and
reproducible benchmark scripts. The implementation aligns the report, demo, and
repository around a single end-to-end workflow: generate candidate reasoning
steps, search over them, score process quality, and return the strongest
trajectory.

\begin{thebibliography}{9}

\bibitem{rstar}
Xinyu Li et al. 2024. rStar-Math: Small LLMs Can Master Math Reasoning with
Self-Evolved Deep Thinking.

\bibitem{mcts}
Levente Kocsis and Csaba Szepesvari. 2006. Bandit Based Monte-Carlo Planning.

\bibitem{math}
Dan Hendrycks et al. 2021. Measuring Mathematical Problem Solving With the MATH
Dataset.

\bibitem{olympiad}
OlympiadBench. A benchmark for olympiad-level bilingual multimodal scientific
reasoning.

\end{thebibliography}

\end{document}
